\section{Objetivos}
\label{sec:objetivos}

\subsection{Objetivo Geral}

Desenvolver e calibrar um modelo compartimental SEIR-SEI, com período de latência nos humanos e o mosquito
\textit{Aedes aegypti} como população explícita, para o monitoramento da chikungunya no Recife, usando as
notificações de 2015 a 2025 e o inquérito sorológico de 2018 e 2019.

\subsection{Objetivos Específicos}

\begin{itemize}
  \item \textbf{OE1} Formular o modelo SEIR-SEI e derivar seus números de reprodução, sua taxa de
  crescimento inicial e a eficácia crítica do controle vetorial.
  \item \textbf{OE2} Organizar a série de casos prováveis de chikungunya por semana e por distrito sanitário
  a partir dos dados abertos da Prefeitura do Recife.
  \item \textbf{OE3} Calibrar o modelo à primeira onda, de 2015 e 2016, estimando o número de reprodução, a
  fração de notificação e a sazonalidade do vetor, com a soroprevalência medida como restrição.
  \item \textbf{OE4} Estender o modelo para grupos da cidade, por distrito sanitário ou por estrato
  socioeconômico, e testar se a heterogeneidade explica a soroprevalência de 36\% e a onda de 2021.
  \item \textbf{OE5} Avaliar cenários de controle vetorial com o modelo calibrado, variando a eficácia e o
  momento de início.
  \item \textbf{OE6} Divulgar os resultados em repositório aberto, relatório técnico e painel de
  monitoramento.
\end{itemize}
